% New writing for the closing chapter of Volume II. This file is never overwritten by split_volumes.py.

Every chapter of this volume ends with its open problems, and this chapter does not restate them. It sorts them by what they would accomplish, so that a reader looking for the next step can see which problems bear on the Riemann Hypothesis itself and which would sharpen a computation. The sorting follows the path of Section~\ref{ov1:sec:path} of Volume I and its three constraints.

\section{Problems on the descent itself}\label{ov2:sec:descent}

These are the problems whose solution would be progress on RH, because they concern the step from level one to the critical line for all heights at once.
\begin{enumerate}
\item What a proof by deconvolution would still need, and why every local computation stops short of it: Section~\ref{cf:sec:gap}.
\item Whether the backward flow is a path at all, and the problems of the level one measure: Sections~\ref{lo:sec:path} and~\ref{lo:sec:open}.
\item The questions on the explicit formula read as a deconvolution, including a rigorous matrix pencil (Problem~\ref{es:prob:rigorous_pencil}): Section~\ref{es:sec:open}.
\item The Balazard, Saias and Yor criterion at every level, and the limits of its Monte Carlo solution: Sections~\ref{bsy:sec:limits} and~\ref{bsy:sec:open}.
\item Where the uniform step of the phase diagram could come from, the de Bruijn Newman flow being the only known source of an analytic tail below the basepoint one (below $\frac78$ if the announced half plane is confirmed): Section~\ref{cp:sec:uniform}.
\end{enumerate}
Each of these meets the first constraint of Section~\ref{ov1:sec:path} of Volume I: a computation in a window cannot settle a statement for all heights (Section~\ref{ch:finite}). A solution has to be an argument, and the Euler product has to enter it.

\section{Problems on windows of the explicit formula}\label{ov2:sec:windows}

Weil positivity in compact windows is certified to finite support lengths, and the barriers of Chapters~\ref{vol2:ch15} and~\ref{vol2:ch16} say how far such certificates can see: Sections~\ref{wt:sec:path}, \ref{wt:sec:open}, \ref{yo:sec:open}, \ref{zh:sec:closer} and~\ref{zh:sec:open}. These problems would extend what is certified. A window certificate is a statement about finitely many support lengths, so the first constraint applies to it as it does to every computation.

\section{Problems on approximants}\label{ov2:sec:approx}

The truncations, the friable approximants, the extreme value approximants and the approximating clocks each agree with $\Xi$ to high accuracy and yet can have zeros off the line. Their open problems ask where the zeros of the approximants go and what a computational criterion must see to tell them from $\Xi$: Sections~\ref{tr:sec:path}, \ref{tr:sec:trnotes}, \ref{mi:sec:open}, \ref{ev:sec:path}, \ref{ev:sec:open}, \ref{ht:sec:need} and~\ref{ht:sec:open}. They test the second and third constraints of Section~\ref{ov1:sec:path} of Volume I: an approximant that keeps the functional equation but not the Euler product, or a property that survives a Gaussian tilt, is where a false criterion would be caught.

\section{Problems on cost and certification}\label{ov2:sec:cost}

The remaining problems concern the computation as such: whether the level one density can be evaluated in time polynomial in both the logarithm of the height and the number of digits (Problem~\ref{cx:prob:polylog}), a posteriori bounds for the accelerated flow, and a calibrated Bayesian inversion of the backward flow (Problem~\ref{lo:prob:sampler}): Sections~\ref{cx:sec:worked} and~\ref{cx:sec:open}. They would make the computations of this volume faster and more rigorous, and they bear on RH only through the windows they make reachable.
